% Partie : applied-modal-logics
% Chapitre : epistemic-logic
% Section : truth-at-w

\documentclass[../../../include/open-logic-section]{subfiles}

\begin{document}

\olfileid[fr]{aml}{el}{trw}

\olsection{Vérité en un monde}

Comme en logique modale normale, chaque modèle épistémique détermine
quelles !!{formula}s sont vraies en quels mondes. Nous employons la
même notation pour la notion fondamentale de la sémantique relationnelle :
« le modèle $\mModel{M}$ rend la !!{formula}~$!A$ vraie au monde~$w$ ».
Cette relation est définie inductivement ; pour tous les opérateurs
non modaux, sa définition est identique à celle de la logique modale normale.

\begin{defn}\ollabel{defn:mmodels}
  La \emph{vérité d'!!a{formula}~$!A$ en~$w$} dans
  un modèle~$\mModel M$, notée $\mSat{M}{!A}[w]$, est définie
  inductivement comme suit :
  \begin{enumerate}
  \tagitem{prvFalse}{\indcase{!A}{\lfalse}{On n'a jamais $\mSat{M}{\lfalse}[w]$}.}{}
  \tagitem{prvTrue}{\indcase{!A}{\ltrue}{On a toujours $\mSat{M}{\ltrue}[w]$}.}{}
  \item $\mSat{M}{p}[w]$ si et seulement si $w \in V(p)$.
  \tagitem{prvNot}{\indcase{!A}{\lnot !B}{$\mSat{M}{\indfrm}[w]$ si et seulement si
    $\mSat/{M}{!B}[w]$}.}{}
  \tagitem{prvAnd}{\indcase{!A}{(!B \land !C)}{$\mSat{M}{\indfrm}[w]$ si et seulement si
    $\mSat{M}{!B}[w]$ et $\mSat{M}{!C}[w]$}.}{}
  \tagitem{prvOr}{\indcase{!A}{(!B \lor !C)}{$\mSat{M}{\indfrm}[w]$ si et seulement si
    $\mSat{M}{!B}[w]$ ou $\mSat{M}{!C}[w]$} (éventuellement les deux).}{}
  \tagitem{prvIf}{\indcase{!A}{(!B \lif !C)}{$\mSat{M}{\indfrm}[w]$ si et seulement si
    $\mSat/{M}{!B}[w]$ ou $\mSat{M}{!C}[w]$}.}{}
  \tagitem{prvIff}{\indcase{!A}{(!B \liff !C)}{$\mSat{M}{\indfrm}[w]$ si et seulement si
    soit on a à la fois $\mSat{M}{!B}[w]$ et $\mSat{M}{!C}[w]$, soit
    on n'a ni $\mSat{M}{!B}[w]$ ni $\mSat{M}{!C}[w]$}.}{}
  \item\ollabel{defn:sub:mmodels-box}
    \indcase{!A}{\Knows_a !B}{$\mSat{M}{\indfrm}[w]$ si et seulement si
    $\mSat{M}{!B}[w']$ pour tout $w' \in W$ tel que $R_a ww'$}
  \end{enumerate}
\end{defn}

C'est ici qu'il faut envisager des restrictions sur nos relations
d'accessibilité. D'après la clause~\olref{defn:sub:mmodels-box},
!!a{formula}~$\Knows_a !B$ est vraie en~$w$ dès qu'il n'existe
aucun~$w'$ tel que $R_a ww'$. Comme en logique modale normale,
toutes les formules en $\Box$ sont vraies par vacuité dans un monde
sans successeur. Ce résultat paraît très contraire à l'intuition si
l'on interprète $\Knows$ comme la \emph{connaissance} : nous tendons
à penser qu'un agent ne peut, en \emph{aucune} circonstance, savoir
à la fois $!A$ et $\lnot !A$.

Une solution consiste à exiger que la relation d'accessibilité soit
\emph{réflexive}. Cela correspond approximativement à l'idée que le
monde réel est compatible avec l'information de l'agent. Les logiques
épistémiques utilisent souvent S5, mais on peut aussi choisir des
systèmes plus faibles selon ce que l'on souhaite représenter par
l'opérateur $\Knows_a$.

\begin{figure}
  \begin{center}
    \begin{tikzpicture}[modal]
      \node[world] (w1) [label={[align=right]right:\mTrue{p}\\ \mFalse{q}}]{$w_1$};
      \node[world] (w2) [label={[align=right]right:\mTrue{p}\\ \mTrue{q}},
        above right=of w1]{$w_2$};
      \node[world] (w3) [label={[align=right]right:\mFalse{p}\\ \mFalse{q}},
        below right=of w1] {$w_3$};
      \draw[<->] (w1) to node {$a$} (w2);
      \draw[->,loop left] (w1) to node {$a,b$} (w1);
      \draw[<->] (w1) to [swap] node {$b$} (w3);
      \draw[->,loop above] (w2) to node {$a,b$} (w2) ;
      \draw[->,loop below] (w3) to node {$a,b$} (w3) ;
    \end{tikzpicture}
  \end{center}
  \caption{Un modèle épistémique simple.}
  \ollabel{fig:simple}
\end{figure}

\begin{prob}
  Déterminez lesquels des énoncés suivants sont vérifiés dans
  la \olref[aml][el][trw]{fig:simple} :
  \begin{enumerate}
    \item $\mSat{M}{\lnot q}[w_1]$;
    \item $\mSat{M}{\Knows_a \lnot q}[w_1]$;
    \item $\mSat{M}{\Knows_b \lnot q}[w_1]$;
    \item $\mSat{M}{\Knows_b q \lor \Knows_b \lnot q}[w_2]$;
    \item $\mSat{M}{\Knows_a( \Knows_b q \lor \Knows_b \lnot q)}[w_2]$;
    \item $\mSat{M}{\EKnows_{\{a,b\}} \lnot q}[w_3]$;
    \end{enumerate}
\end{prob}

Une fois la vérité en un monde définie, les autres notions sémantiques
de la logique modale normale, telles que la validité modale et la
conséquence logique, se transposent à cette nouvelle interprétation
des opérateurs modaux.

Nous pouvons maintenant donner les conditions de vérité de l'opérateur
de connaissance commune~$\CKnows_G$. Rappelons, d'après la
\olref[sfr][rel][ops]{sec}, que la \emph{clôture transitive}~$R^+$
d'une relation~$R$ est définie comme suit.
\footnote{L'indice zéro ci-dessous correspond à un pas : l'indexation
est décalée d'une unité par rapport à celle de la section citée.
La relation obtenue reste la clôture des chemins de longueur
strictement positive.}
\begin{align*}
  R^+ &= \bigcup_{ n \in \mathbb{N}} R^n,
\intertext{où}
  R^0 & = R \text{ et}\\
  R^{n+1} & = \Setabs{\tuple{x, z}}{\exists y (R^n xy \land Ryz)}.
\end{align*}
Pour un groupe d'agents $G$, on définit alors
$R_G = ( \bigcup_{b
\in G} R_b )^+$ comme la clôture transitive de la réunion des relations
d'accessibilité de ses agents.

\begin{defn}
Si $G' \subseteq G$, on pose $\mSat{M}{\CKnows_{G'} !A}[ w]$ si et
seulement si, pour tout $w'$ tel que $R_{G'} w w'$, on a
$\mSat{M}{!A}[w']$.
\end{defn}

Cette définition retient toutes les itérations de connaissance de
longueur strictement positive. Si les relations des agents sont
réflexives et le groupe non vide, elle inclut aussi la vérité de la
formule au monde de départ. Sur des relations arbitraires, elle
n'ajoute pas automatiquement cette condition.

\end{document}
